% ==============================================================================
% 《形形色色的曲线》小 32 开 (130mm x 184mm) 忠实重排版统一样式配置
% ==============================================================================
\usepackage[
  paperwidth=130mm,
  paperheight=184mm,
  top=15.5mm,
  bottom=14.5mm,
  left=14mm,
  right=14mm,
  headheight=20pt,
  headsep=4mm,
  footskip=6.5mm
]{geometry}

% 数学与字体宏包
\usepackage{amsmath, amssymb}
\usepackage{newtxmath}
\let\widering\relax
\usepackage{yhmath} % 提供平滑圆弧符号 \wideparen

\setmainfont{TeX Gyre Termes}

% TikZ 与 PGFPlots 矢量绘图
\usepackage{tikz}
\usepackage{pgfplots}
\pgfplotsset{compat=1.18}
\usetikzlibrary{
  arrows.meta,
  decorations.pathreplacing,
  decorations.markings,
  angles,
  quotes,
  calc,
  intersections,
  patterns
}

% 图表与混排控制
\usepackage{graphicx}
\graphicspath{{assets/generated/}{assets/figures/}}
\usepackage{wrapfig}
\usepackage{caption}
\captionsetup{
  font=footnotesize,
  labelfont=bf,
  labelsep=quad,
  skip=4pt
}
\usepackage{longtable, booktabs, array}

% 颜色、页眉页脚、标题与目录
\usepackage{xcolor}
\usepackage{fancyhdr}
\usepackage{titlesec}
\usepackage{titletoc}
\usepackage[
  colorlinks=true,
  linkcolor=accentblue,
  urlcolor=accentblue,
  bookmarksnumbered=true,
  pdfstartview=FitH
]{hyperref}

% 配色方案
\definecolor{accentblue}{RGB}{24, 72, 122}
\definecolor{curveblue}{RGB}{20, 90, 160}
\definecolor{mutedgray}{RGB}{95, 100, 110}
\definecolor{coverbg}{RGB}{66, 92, 82}
\definecolor{covergreen}{RGB}{118, 224, 128}
\definecolor{coverdark}{RGB}{44, 62, 55}

% ------------------------------------------------------------------------------
% 原书数学符号忠实还原宏定义
% ------------------------------------------------------------------------------
\DeclareMathOperator{\tg}{tg}
\DeclareMathOperator{\ctg}{ctg}
\DeclareMathOperator{\arctg}{arctg}
\DeclareMathOperator{\arcctg}{arcctg}
\DeclareMathOperator{\sh}{sh}
\DeclareMathOperator{\ch}{ch}
\DeclareMathOperator{\thyp}{th}
\DeclareMathOperator{\cth}{cth}
\newcommand{\arc}[1]{\wideparen{#1}}
\newcommand{\dd}{\,d}

% ------------------------------------------------------------------------------
% 节标题与目录样式（扁平 100 节结构）
% ------------------------------------------------------------------------------
\renewcommand{\thesection}{\arabic{section}}
\titleformat{\section}
  {\centering\large\sffamily\bfseries\color{accentblue}}
  {\thesection.}{0.45em}{}
\titlespacing*{\section}{0pt}{1.0\baselineskip}{0.65\baselineskip}

% 紧凑清晰的目录格式（小 32 开单列对齐，纯文本 contentsname 避免样式污染页眉）
\setcounter{tocdepth}{1}
\renewcommand{\contentsname}{目录}
\ctexset{
  chapter = {
    name = {},
    number = {},
    format = \centering\Large\sffamily\bfseries\color{accentblue},
    beforeskip = 0.8\baselineskip,
    afterskip = 0.8\baselineskip,
    pagestyle = fancy
  }
}
\titlecontents{section}
  [0em]
  {\small}
  {\makebox[2.2em][r]{\thecontentslabel.\hspace{0.4em}}}
  {}
  {\titlerule*[0.5pc]{.}\contentspage}

% ------------------------------------------------------------------------------
% 页眉与页脚设置
% ------------------------------------------------------------------------------
\pagestyle{fancy}
\fancyhf{}
\fancyhead[LE,RO]{\footnotesize\color{mutedgray} 形形色色的曲线}
\fancyhead[LO,RE]{\footnotesize\color{mutedgray} \leftmark}
\fancyfoot[C]{\footnotesize\thepage}
\renewcommand{\headrulewidth}{0.4pt}
\renewcommand{\headrule}{\hbox to\headwidth{\color{mutedgray!40}\leaders\hrule height \headrulewidth\hfill}}

% 允许较长多行公式适当跨页，减少大段空白
\allowdisplaybreaks[2]

% 适配小 32 开窄版心的正文字距与行距（小五号约 9.5pt，保证公式不溢出）
\raggedbottom
\linespread{1.28}
\setlength{\parindent}{2em}
\setlength{\parskip}{0.15\baselineskip}
